\subsection{Representation-Adaptive Side Information}

The fixed-length theorem pays for the worst collision fiber through $A_\pi$. A finer coding law allows the auxiliary description length to depend on the observed representation value $u$. Then the relevant budget is not the maximum fiber size but the size of the individual fiber $\pi^{-1}(u)$. The same residual ambiguity object is now resolved fiber by fiber and then averaged under the induced law of $U$.

\begin{definition}[Representation-adaptive auxiliary description]
An \emph{adaptive auxiliary scheme} assigns to each representation value $u$ a binary description budget $\ell(u)$ and a decoder that, given $u$ together with a binary string of length at most $\ell(u)$, recovers the class exactly whenever $\pi(C)=u$.
\end{definition}

\begin{theorem}[Pointwise optimal adaptive bit budget]\label{thm:adaptive-fiber-budget}
For a representation fiber of size $|\pi^{-1}(u)|$, zero-error recovery under a binary auxiliary alphabet is feasible if and only if
\[
|\pi^{-1}(u)| \le 2^{\ell(u)}.
\]
Consequently, the pointwise optimal number of auxiliary bits on fiber $u$ is
\[
\ell^\star(u) = \left\lceil \log_2 |\pi^{-1}(u)| \right\rceil.
\]
\end{theorem}
\leanmeta{\LH{GPH32}}

\begin{proof}
Once $u$ is fixed, the decoder only needs to distinguish classes inside the fiber $\pi^{-1}(u)$. With $\ell(u)$ binary bits there are at most $2^{\ell(u)}$ auxiliary outcomes, so exact recovery is possible exactly when that many outcomes suffice to name the fiber elements injectively. The minimal feasible budget is therefore the least integer whose power of two dominates the fiber size, namely $\lceil \log_2 |\pi^{-1}(u)| \rceil$.
\end{proof}

\begin{definition}[Expected adaptive description length]
Let $U=\pi(C)$ be the induced representation random variable under a source law on classes. For an adaptive budget $\ell$, define its expected auxiliary length by
\[
\mathbb E[\ell(U)] = \sum_u P_U(u)\,\ell(u).
\]
\end{definition}

\begin{theorem}[Adaptive expected-rate lower bound]\label{thm:adaptive-expected-rate-lower-bound}
Any feasible adaptive auxiliary scheme satisfies
\[
\mathbb E[\ell(U)] \ge \mathbb E\!\left[\left\lceil \log_2 |\pi^{-1}(U)| \right\rceil\right].
\]
Moreover, the pointwise optimal budget from Theorem~\ref{thm:adaptive-fiber-budget} attains equality.
\end{theorem}
\leanmeta{\LHrng{GPH}{32}{33}}

\begin{proof}
By Theorem~\ref{thm:adaptive-fiber-budget}, feasibility on fiber $u$ requires
\[
\ell(u) \ge \left\lceil \log_2 |\pi^{-1}(u)| \right\rceil
\]
for every representation value $u$. Multiplying pointwise by $P_U(u)$ and summing over $u$ gives the lower bound on expected length. Equality is achieved by choosing the pointwise optimal budget on every fiber.
\end{proof}

\begin{corollary}[Conditional-entropy lower bound]
Let $C$ be drawn from any source distribution and let $U=\pi(C)$. Then any feasible adaptive auxiliary scheme satisfies
\[
\mathbb E[\ell(U)] \ge \frac{H(C\mid U)}{\log 2},
\]
where $H(C\mid U)$ is conditional Shannon entropy in natural-log units.
\end{corollary}
\leanmeta{\LH{GPH34}}

\begin{proof}
For each representation value $u$, the conditional law of $C$ given $U=u$ is supported on the fiber $\pi^{-1}(u)$, so its entropy is at most $\log |\pi^{-1}(u)|$. Averaging over $u$ gives
\[
H(C\mid U) \le \mathbb E[\log |\pi^{-1}(U)|].
\]
Since $\log |\pi^{-1}(u)| \le \lceil \log_2 |\pi^{-1}(u)| \rceil \log 2$ fiberwise, Theorem~\ref{thm:adaptive-expected-rate-lower-bound} yields the stated lower bound.
\end{proof}

\begin{theorem}[Fiberwise prefix-coding upper bound]\label{thm:fiberwise-prefix-coding-upper-bound}
Assume standard prefix-coding achievability on each finite fiber. Then there exists an adaptive zero-error auxiliary code satisfying
\[
\mathbb E[\ell(U)] \le \frac{H(C\mid U)}{\log 2} + 1.
\]
\end{theorem}
\leanmeta{\LH{GPH35}}

\begin{proof}
Because the decoder already knows $U=u$, prefix-freeness is required only within each fiber $\pi^{-1}(u)$. Applying a standard prefix-coding theorem separately to the conditional source on each nonempty fiber gives an expected conditional code length at most $H(C\mid U=u)/\log 2 + 1$. Averaging over $u$ yields the stated bound.
\end{proof}

\begin{corollary}[Zero-error conditional-entropy sandwich]\label{cor:zero-error-conditional-entropy-sandwich}
The optimal expected zero-error auxiliary rate lies between
\[
\frac{H(C\mid U)}{\log 2}
\qquad\text{and}\qquad
\frac{H(C\mid U)}{\log 2}+1.
\]
\end{corollary}
\leanmeta{\LH{ZEC1}}

\begin{proof}
The lower bound is Theorem~\ref{thm:adaptive-expected-rate-lower-bound}, and the upper bound is Theorem~\ref{thm:fiberwise-prefix-coding-upper-bound}. The two statements together give the claimed sandwich.
\end{proof}

\begin{corollary}[Active binary clarification]
If, after observing $U$, the decoder is allowed to ask adaptive binary clarification questions whose leaves identify the class exactly, then the minimum expected number of clarification questions is bounded between
\[
\frac{H(C\mid U)}{\log 2}
\qquad\text{and}\qquad
\frac{H(C\mid U)}{\log 2}+1.
\]
\end{corollary}
\leanmeta{\LHrng{GPH}{34}{35}}

\begin{proof}
Each binary clarification protocol induces a binary code over the conditional source on each fiber, and conversely any fiberwise binary prefix code can be implemented as an adaptive yes/no decision tree. The lower and upper bounds therefore follow from the preceding conditional-entropy comparison and fiberwise prefix-coding theorem.
\end{proof}

\begin{corollary}[Representation-respecting noisy transcripts cannot beat the fiber law]
Let $T$ be any repeated or noisy observation transcript whose realized value factors through the fixed representation, that is, $T = h(U)$ for some map $h$. Then recovering $C$ from $T$ cannot require less worst-case or pointwise adaptive side information than recovering $C$ from $U$ itself.
\end{corollary}
\leanmeta{\LH{GPH28}, \LH{GPH37}}

\begin{proof}
If the entire transcript factors through $U$, then it is a coarsening of the original representation. Corollary~\ref{cor:concept-bottleneck-mono} applies directly: coarsening cannot reduce the largest collision fiber and cannot reduce the pointwise adaptive fiber budget. Thus repeated or noisy observations that never break the original fibers cannot remove the residual zero-error coding burden.
\end{proof}

\begin{remark}[Bit currency at fiber resolution]
The fixed-length law pays for the worst collision fiber; the adaptive law pays for each fiber separately and averages under the source-induced law of $U$. These are two resolutions of the same object: the cost of naming the classes that remain collapsed inside each representation fiber.
\end{remark}

\begin{remark}[Entropy comparison]
The conditional-entropy lower and upper bounds place the adaptive theorem inside classical source-coding language. Zero-error adaptive side information is sandwiched around $H(C\mid U)$ up to the standard one-bit prefix-coding overhead, but the exact zero-error law remains fiberwise and combinatorial rather than purely entropic.
\end{remark}
